% Smaller copies of the message diagram from slides/02-cvqc.tex.
% Arguments: vertical position, verifier label.
\newcommand{\strategyprotocol}[2]{%
  \begin{scope}[shift={(0,#1)},scale=0.68,transform shape]
    \node[party] at (-4.15,0) {#2};
    \node[party] at (4.15,0) {Prover};

    \draw[communication] (-2.35, 1.10) -- ( 2.35, 1.10);
    \draw[communication] ( 2.35, 0.50) -- (-2.35, 0.50);
    \draw[communication] (-2.35,-0.10) -- ( 2.35,-0.10);
    \node[font=\large,text=CVQCTeal] at (0,-0.60) {$\vdots$};
    \draw[communication] ( 2.35,-1.10) -- (-2.35,-1.10);
  \end{scope}%
}

% outline.md, lines 37--45.
\begin{frame}{CVQC Strategy}
  \centering
  \begin{tikzpicture}[
      >={Latex[length=3mm]},
      party/.style={circle,draw=CVQCInk,fill=CVQCFog,very thick,
        minimum size=2.5cm,text width=2.1cm,inner sep=0pt,
        align=center,font=\large\bfseries,text=CVQCInk},
      communication/.style={->,line width=1.5pt,draw=CVQCTeal},
      protocolbox/.style={draw=CVQCInk,thick,rounded corners=4mm},
      modellabel/.style={fill=white,inner xsep=6pt,inner ysep=2pt,
        font=\bfseries,text=CVQCInk}
    ]
    % Keep the first diagram in place when revealing the compiler below it.
    \path[use as bounding box] (-6.5,-3.15) rectangle (6.5,3.15);

    \draw[protocolbox] (-5.4,0.4) rectangle (5.4,2.9);
    \node[modellabel] at (-4.25,2.9) {$\mathsf{QPIP}_1$};
    \strategyprotocol{1.65}{Weak\\quantum\\verifier}

    \only<2>{
      \draw[->,line width=1.5pt,draw=CVQCTeal] (0,0.30) -- (0,-0.30)
        node[midway,right=8pt,font=\small\bfseries,text=CVQCTeal] {compile};

      \draw[protocolbox] (-5.4,-2.9) rectangle (5.4,-0.4);
      \node[modellabel] at (-4.25,-0.4) {$\mathsf{QPIP}_0$};
      \strategyprotocol{-1.65}{Classical\\verifier}
    }
  \end{tikzpicture}

  \note<1>{
    \footnotesize
    A common strategy is to begin with a verifier that has a small amount of
    quantum capability, interacting with a full quantum prover. We refer to
    this model as $\mathsf{QPIP}_1$. In our construction, the verifier only
    prepares and sends individual qubits from a fixed finite set; the rest of
    the interaction is classical. The arrows indicate interaction, not a
    particular number of messages.
  }
  \note<2>{
    \footnotesize
    Next, compile the protocol to remove the verifier's quantum operations,
    giving a fully classical verifier interacting with the quantum prover:
    the $\mathsf{QPIP}_0$ model. The bottom protocol is the compiled version
    of the one above. This is the overall strategy; the next slides show
    two paradigms for implementing this compilation. In our prepare-and-send setting, the
    quantum preparation is replaced with remote state preparation.
  }
\end{frame}
